\documentclass[10pt,a4paper]{amsart}
\usepackage[margin=19mm]{geometry}
\usepackage{amsmath,amssymb,mathtools}
\usepackage[scr=rsfs,cal=euler]{mathalfa}
\usepackage{xfrac}
\usepackage[unicode,hidelinks,bookmarks=false]{hyperref}
\newcommand{\ie}{i.e.\ }
\newcommand{\cf}{cf.\ }
\newcommand{\resp}{respectively\ }
\newcommand{\Subsection}{\subsection}
\providecommand{\nobreakdash}{\nobreak\hbox{-}}
\setlength{\emergencystretch}{2em}
\multlinegap0pt
\usepackage{bm}
\usepackage{bbm}
\usepackage{dsfont}
\usepackage{xspace}
\usepackage{cancel}
\usepackage{mathtools}
\usepackage{amsthm}
\makeatletter
\@ifundefined{NB}{\newtheorem{NB}{Nb}}{}
\@ifundefined{thm}{\newtheorem{thm}{Theorem}}{}
\@ifundefined{lem}{\newtheorem{lem}[thm]{Lemma}}{}
\@ifundefined{prop}{\newtheorem{prop}[thm]{Proposition}}{}
\makeatother
\newcommand{\NS}{\operatorname{NS}}
\newcommand{\cal}{\mathcal}
\title[Preservation of stability under the Fourier-Mukai transform ]{Preservation of stability under the Fourier-Mukai transform whose kernel is the Poincare line bundle}
\author{Kota Yoshioka}
\date{Source: arXiv:2506.17567v2 (CC BY 4.0)}
\begin{document}
\maketitle

\noindent\emph{Source and licence.} Preservation of stability under the Fourier-Mukai transform whose kernel is the Poincare line bundle. Version arXiv:2506.17567v2 (CC BY 4.0), distributed under the Creative Commons Attribution 4.0 International licence, as recorded in the arXiv licence metadata for this exact version and shown on the abstract page https://arxiv.org/abs/2506.17567v2. Full rights notice: licenses/arxiv-2506.17567-CC-BY-4.0.txt.\par\smallskip

\noindent\textit{Editorial scope.} Counted unit: the Thm whose source label is \texttt{thm:stability} in the article (source0.tex lines 1719--1743, source label \texttt{thm:stability}), delivered together with the article's own complete original proof of it (source0.tex lines 1745--1856, one \texttt{proof} environment of 112 lines). No mathematics is written, repaired or re-derived: statement and proof are the article's own text line for line. Required context retained uncounted: prop:isom (lines 1066--1085); eq:u (lines 1127--1130); prop:t\_0 (lines 1298--1312); lem:Gieseker (lines 1399--1403); lem:a>0 (lines 1511--1551); lem:a<0 (lines 1632--1664). Every definition, display and label of those lines is retained unchanged. Omitted from this package and not redistributed: 1--1065, 1086--1126, 1131--1297, 1313--1398, 1404--1510, 1552--1631, 1665--1718, 1744--1744, 1857--3160. Nothing inside a retained excerpt is omitted or altered: every retained body is delivered exactly as the article's lines are written, so no omission receipt of any kind accompanies this package. Citations: the retained excerpts carry no citation command at all, so no bibliography entry of the article is transcribed and no reference is redistributed. Reference adaptation: none was needed and none was made; every reference inside the delivered unit resolves inside it, so no unresolved reference or citation is produced. Printed slips kept literal: no printed word, symbol, hypothesis or number of the counted statement or of its proof was altered. Layout and class: the article's own class is \texttt{amsart} with packages including amssymb, amsbsy, amscd, eucal, verbatim, version, color, xy, graphicx; the wrapper is the lane's pinned \texttt{amsart} editorial wrapper at 10pt on A4, which restates the theorem-like environments the retained text needs and adds the article's own macro definitions that the retained text needs (\texttt{\textbackslash NS}, \texttt{\textbackslash cal}), with extra wrapper packages (bm, bbm, dsfont, xspace, cancel, mathtools). The delivered numbering is the unit's own. Assets: no retained span contains a figure environment, a graphics-inclusion command, a PDF-page inclusion, a table or a TikZ picture; no image file is copied into this package, the package declares no asset dependency and no image file is redistributed with it. Licence: version arXiv:2506.17567v2 of this article is distributed under the Creative Commons Attribution 4.0 International licence, as recorded in the arXiv licence metadata for this exact version and shown on the abstract page https://arxiv.org/abs/2506.17567v2. A full rights notice is delivered at licenses/arxiv-2506.17567-CC-BY-4.0.txt. Characters: every retained excerpt is pure ASCII, so the delivered engine, XeLaTeX with its default Latin Modern text font, reports no missing character.
\begin{prop}\label{prop:isom}
\begin{enumerate}
\item[(1)]
$\Phi[1]$ induces an isomorphism
$$
{\cal M}_{(0,tH)}(r,\xi,a) \to {\cal M}_{(0,(nt)^{-1} \widehat{H})}(-a,\widehat{\xi},-r).
$$
\item[(2)]
Assume that $t \ll 1$.
\begin{enumerate}
\item
If $a \leq 0$, then
${\cal M}_{(0,(nt)^{-1} \widehat{H})}(-a,\widehat{\xi},-r)={\cal M}_{\widehat{H}}(-a,\widehat{\xi},-r)$.
\item
If $a>0$, then
${\cal M}_{(0,(nt)^{-1} \widehat{H})}(-a,\widehat{\xi},-r)$ 
consists of $E^{\vee}[1]$, $E \in {\cal M}_{\widehat{H}}(a,-\widehat{\xi},r)$.
\end{enumerate}
\end{enumerate}
\end{prop}

\begin{equation}\label{eq:u}
v=\ell_1 v_1+\ell_2 v_2, \; \langle v_1^2 \rangle=\langle v_2^2 \rangle=0,\; 
\langle v_1,v_2 \rangle=1,\; \{ \ell_1,\ell_2 \}=\{ \ell,1 \},\; u \in \{v_1,v_2 \}. 
\end{equation}

\begin{prop}\label{prop:t_0}
There are real numbers $t_1 \geq t_2 \geq 0$ such that
\begin{enumerate}
\item
if $t>t_1$, then 
$E$ is a torsion free sheaf for a general $E \in {\cal M}_{(0,tH)}(v)$ and
\item
if $t_1>t>t_2$, then
$E$ is a coherent sheaf with torsions for a general
$E \in {\cal M}_{(0,tH)}(v)$ and
\item
if $t<t_2$, then
 $H^{-1}(E) \ne 0$ for a general $E \in {\cal M}_{(0,tH)}(v)$.
\end{enumerate}
\end{prop}

\begin{lem}\label{lem:Gieseker}
Assume that $t_1+\epsilon>t>t_1>0$, where $\epsilon$ is sufficiently small positive number.
Then there is an ample divisor $L$ such that
${\cal M}_{(0,tH)}(v) \cap {\cal M}_{L}(v) \ne \emptyset$.
\end{lem}

\begin{lem}\label{lem:a>0}
Assume that $a >0$.
Then $a_1>0$,
$\Phi(E_1) \in
{\cal M}_{\widehat{H}}(\ell_1(a_1,-\widehat{\xi_1},r_1))$
and the following claims hold.
\begin{enumerate}
\item[(1)]
Assume that $r_2<0$.
Then $a_2>0$ and $\Phi(E_2) \in {\cal M}_{\widehat{H}}(\ell_2(a_2,-\widehat{\xi_2},r_2))$.
%or $v_2=(-1,0,0)$ and $\Phi(E_2)[1]$ is a 0-dimensional sheaf.
\item[(2)]
Assume that $r_2=0$.
Then one of the following holds.
\begin{enumerate}
\item
$a_2>0$ and $\Phi(E_2) \in {\cal M}_{\widehat{H}}(\ell_2(a_2,-\widehat{\xi_2},r_2))$.
\item
$a_2=0$, 
$(\xi_2^2)=0$, $(\xi_1 \cdot \xi_2)=1$
and $\Phi(E_2)[1] \in {\cal M}_{\widehat{H}}(\ell_2(0,\widehat{\xi_2},0))$.
In particular $X$ is a product of elliptic curves and $\xi$ is primitive.
\begin{NB}
$(\xi \cdot \xi_i)=\ell_{1-i}$.
Hence $\xi$ is primitive. 
\end{NB}
\end{enumerate}
\item[(3)]
Assume that $r_2>0$. Then the following claims hold.
\begin{enumerate}
\item
If $a_2 > 0$, then
$\Phi(E_2) \in {\cal M}_{\widehat{H}}(\ell_2(a_2,-\widehat{\xi_2},r_2))$.
\item
If $a_2 \leq 0$, then 
%$\Phi(E_2)[1] \in {\cal M}_{\widehat{H}}(0,\widehat{\xi_2},-r_2)$ is a torsion sheaf.
%If $a_2<0$, then
$\Phi(E_2)[1] \in {\cal M}_{\widehat{H}}(\ell_2(-a_2,\widehat{\xi_2},-r_2))$. 
\end{enumerate}
\end{enumerate}
\end{lem}

\begin{lem}\label{lem:a<0}
Assume that $a \leq 0$.
Then
$a_2<0$,
$\Phi(E_2)[1] \in {\cal M}_{\widehat{H}}(\ell_2(-a_2,\widehat{\xi_2},-r_2))$
and the following claims hold.
\begin{enumerate}
\item[(1)]
Assume that $r_2<0$.
Then $a_1<0$ and $\Phi(E_1)[1] \in 
{\cal M}_{\widehat{H}}(\ell_1(-a_1,\widehat{\xi_1},-r_1))$.
\item[(2)]
Assume that $r_2=0$. Then one of the following holds.
\begin{enumerate}
\item
If $a_1=0$, then there is an elliptic curve $C$ and
$v_1=(1,k_1 C,0)$, $v_2=(0,k_2 C,-1)$, where $k_1,k_2 \in {\Bbb Z}_{>0}$.
\item
If $a_1<0$, then $\Phi(E_1)[1] \in {\cal M}_{\widehat{H}}(\ell_1(-a_1,\widehat{\xi_1},-r_1))$.
\end{enumerate}
\item[(3)]
\begin{NB}
This case does not occur for the proof of Theorem \ref{thm:stability} (2).
\end{NB}
Assume that $r_2>0$.
\begin{enumerate}
\item
 If $a_1 \leq 0$, then $\Phi(E_1)[1] \in {\cal M}_{\widehat{H}}(\ell_1(-a_1,\widehat{\xi_1},-r_1))$.
\item
If $a_1>0$, then $\Phi(E_1) \in  {\cal M}_{\widehat{H}}(\ell_1(a_1,-\widehat{\xi_1},r_1))$.
\end{enumerate}
\end{enumerate}
\end{lem}

\begin{thm}\label{thm:stability}
Let $v=(r,\xi,a)$ be a Mukai vector such that $r > 0$ and $(\xi \cdot H)>0$, where 
$H$ is an ample divisor on $X$.
%$v \ne ((\xi_1^2)/2-\ell,\xi_1,1),(\ell (\xi_1^2)/2-1,\ell \xi_1,\ell)$ and
\begin{enumerate}
\item[(1)]
Assume that $a>0$. 
If $X$ is not a product of elliptic curves or $\xi$ is not primitive,
then there are ample divisors $L \in \NS(X)$ and $L' \in \NS(\widehat{X})$ such that 
$\Phi(E)^{\vee}$ is stable with respect to $L'$ for a general
$E \in {\cal M}_L(v)$.
\item[(2)]
Assume that $a \leq 0$ and
$v \ne (\ell,kC,-1), (1,kC,-\ell)$, where 
$C$ is an elliptic curve and $k \geq \ell+1$. 
Then there are ample divisors $L \in \NS(X)$ and $L' \in \NS(\widehat{X})$ such that 
$\Phi(E)[1]$ is stable with respect to $L'$ for a general
$E \in {\cal M}_L(v)$.
\end{enumerate}
%for some $H$ and $H'$,
%where $w=v(\Phi(E))$ or $v=(\Phi(E)[1])$ according as 
%$a>0$ or $a \leq 0$.
%there is a chamber ${\cal C}$ such that $\Phi(E)$ is a stable sheaf up to shift.
%with respect to an ample divisor $H'$.
\end{thm}

\begin{proof}
\begin{NB}
We may assume that $H$ is a general ample divisor with respect to $v$.
\end{NB}
For the family of stability conditions
$\sigma_{(0,tH)}$, let $t_1 \geq t_2$ be non-negative numbers in Proposition \ref{prop:t_0}.
If $t_1+\epsilon> t>t_1>0$, then Lemma \ref{lem:Gieseker} implies
${\cal M}_{(0,tH)}(v) \cap {\cal M}_L(v) \ne \emptyset$ for an ample divisor $L$.

(1)
Assume that $t_1>t_2$ and take a member $E \in {\cal M}_{(0,tH)}(v)$ ($t_1+\epsilon> t>t_1$).
Let   
\begin{equation}\label{eq:HNFt1}
E_1 \to E \to E_2 \to E_1[1]
\end{equation}
be the exact triangle
which is the Harder-Narasimhan filtration of $E$ 
with respect to $\sigma_{(0,tH)}$-semistability $(t_1-\epsilon<t<t_1)$. 
Then $v(E_1)=\ell_1 v_1$ and $v(E_2)=\ell_2 v_2$ in the notation of \eqref{eq:u}.
By the proof of Proposition \ref{prop:t_0},
$E_2$ is a torsion semi-homogeneous sheaf.
By Lemma \ref{lem:a>0}, $\Phi(E_1)$ is a semi-homogeneous bundle.
Since $X$ is not a product of elliptic curves or $\xi$ is not primitive,
$\Phi(E_2)$
%\in \Coh(\widehat{X})$
is also a semi-homogeneous bundle (Lemma \ref{lem:a>0} (2)).
Hence we have an exact sequence
$$
0 \to \Phi(E_2)^{\vee} \to \Phi(E)^{\vee} \to \Phi(E_1)^{\vee} \to 0.
$$
By Lemma \ref{lem:Gieseker}, 
$\Phi(E)^{\vee} \in {\cal M}_{L'}(a,\widehat{\xi},r)$ for an ample divisor $L'$ 
on $\widehat{X}$.

Assume that $t_1=t_2>0$
and take a member $E \in {\cal M}_{(0,tH)}(v)$ ($t_1+\epsilon> t>t_1$).
Let   
\begin{equation}\label{eq:HNFt0}
E_1 \to E \to E_2 \to E_1[1]
\end{equation}
be the exact triangle
which is the Harder-Narasimhan filtration of $E$ 
with respect to $\sigma_{(0,tH)}$-semistability $(t_2-\epsilon <t<t_2)$. 
Then $v(E_1)=\ell_1 v_1$ and $v(E_2)=\ell_2 v_2$ in the notation of \eqref{eq:u}.
By the proof of Proposition \ref{prop:t_0},
$E_2[-1]$ is a semi-homogeneous vector bundle.
%If $\rk \Phi(E_2)=0$, then $\Phi(E)^{\vee}$ is torsion free
By Lemma \ref{lem:a>0},
%$\rk \Phi(E_2) > 0$. 
$\Phi(E_1) $ and $\Phi(E_2)$
%\in \Coh(\widehat{X})$
are semi-homogeneous vector bundles
and we have an exact sequence
$$
0 \to \Phi(E_2)^{\vee} \to \Phi(E)^{\vee} \to \Phi(E_1)^{\vee} \to 0.
$$
By Lemma \ref{lem:Gieseker}, 
$\Phi(E)^{\vee} \in {\cal M}_{L'}(a,\widehat{\xi},r)$ for an ample divisor $L'$
on $\widehat{X}$.

Assume that $t_1=t_2=0$.
We take a general $E \in {\cal M}_H(v)$ such that  
$E \in  {\cal M}_{(0,tH)}(v)$ ($\epsilon> t>0$).
Then  $\Phi(E)^{\vee}$ is a stable sheaf with respect to $\widehat{H}$
by Proposition \ref{prop:isom}.
%For a general $E \in  {\cal M}_{(0,tH)}(v)$ ($\epsilon> t>0$),
%$\Phi(E)^{\vee}$ is a stable sheaf with respect to $\widehat{H}$.

(2)
Assume that $t_1>t_2$ and take a member $E \in {\cal M}_{(0,tH)}(v)$ ($t_1+\epsilon> t>t_1$)
fitting in an exact triangle
$$
E_1 \to E \to E_2 \to E_1[1]
$$
as in \eqref{eq:HNFt1}.
Since $E_2$ is a torsion semi-homogeneous sheaf,
by using Lemma \ref{lem:a<0},
we see that $\Phi(E_1)[1]$ and $\Phi(E_2)[1]$
%\in \Coh(\widehat{X})$
are semi-homogeneous bundles
fitting in an exact sequence
$$
0 \to \Phi(E_1)[1] \to \Phi(E)[1] \to \Phi(E_2)[1] \to 0.
$$
By Lemma \ref{lem:Gieseker}, 
$\Phi(E)[1] \in {\cal M}_{L'}(-a,\widehat{\xi},-r)$ for an ample divisor $L'$
on $\widehat{X}$.

Assume that $t_1=t_2>0$
and take a member $E \in {\cal M}_{(0,tH)}(v)$ ($t_1+\epsilon> t>t_1$)
fitting in an exact triangle
$$
E_1 \to E \to E_2 \to E_1[1]
$$
as in \eqref{eq:HNFt0}.
Since $r_2<0$, by using Lemma \ref{lem:a<0},
we see that $\Phi(E_1)[1] $ and $\Phi(E_2)[1]$
%\in \Coh(\widehat{X})$
are semi-homogeneous bundles
and we have an exact sequence
$$
0 \to \Phi(E_1)[1] \to \Phi(E)[1] \to \Phi(E_2)[1] \to 0.
$$
By Lemma \ref{lem:Gieseker}, 
$\Phi(E)[1] \in {\cal M}_{L'}(-a,\widehat{\xi},-r)$ for an ample divisor $L'$
on $\widehat{X}$.

Assume that $t_1=t_2=0$. We take a general $E \in {\cal M}_H(v)$ such that  
$E \in  {\cal M}_{(0,tH)}(v)$ ($\epsilon> t>0$).
Then  $\Phi(E)[1]$ is a stable sheaf with respect to $\widehat{H}$
by Proposition \ref{prop:isom}.
\end{proof}

\end{document}
